\documentclass[../../main.tex]{subfiles}
\begin{document}
\section{Hamilton 向量场}\label{section:ch22-hamiltonian-vector-fields}

辛流形上最有用的构造之一，是梯度的一个辛类似物，定义如下。设 $(M,\omega)$ 是辛流形。对任意光滑函数 $f\in C^\infty(M)$，我们把 $f$ 的\emph{Hamilton 向量场}定义为由下式给出的光滑向量场 $X_f$
\begin{align*}
X_f=\widetilde\omega^{-1}(\mathrm{d}f),
\end{align*}
其中 $\widetilde\omega:TM\to T^*M$ 是由 $\omega$ 确定的丛同构。等价地，
\begin{align*}
X_f\lrcorner\,\omega=\mathrm{d}f,
\end{align*}
或对任意向量场 $Y$，
\begin{align*}
\omega(X_f,Y)=\mathrm{d}f(Y)=Yf.
\end{align*}

在任意 Darboux 坐标下，$X_f$ 可显式计算如下。写 $X_f=\sum_{i=1}^n\big(a^i\,\partial/\partial x^i+b^i\,\partial/\partial y^i\big)$，其中系数函数 $(a^i,b^i)$ 待定。我们计算
\begin{align*}
X_f\lrcorner\,\omega
=\sum_{j=1}^n\big(a^j\partial/\partial x^j+b^j\partial/\partial y^j\big)\lrcorner\sum_{i=1}^n\mathrm{d}x^i\wedge\mathrm{d}y^i
=\sum_{i=1}^n\big(a^i\,\mathrm{d}y^i-b^i\,\mathrm{d}x^i\big).
\end{align*}
另一方面，
\begin{align*}
\mathrm{d}f=\sum_{i=1}^n\big(\tfrac{\partial f}{\partial x^i}\,\mathrm{d}x^i+\tfrac{\partial f}{\partial y^i}\,\mathrm{d}y^i\big).
\end{align*}
令这两个表达式相等，我们得到 $a^i=\partial f/\partial y^i$ 和 $b^i=-\partial f/\partial x^i$，从而得到 Darboux 坐标下 $f$ 的 Hamilton 向量场公式
\begin{align}
X_f=\sum_{i=1}^n\big(\tfrac{\partial f}{\partial y^i}\,\tfrac{\partial}{\partial x^i}
-\tfrac{\partial f}{\partial x^i}\,\tfrac{\partial}{\partial y^i}\big).
\end{align}
这个公式特别地在带标准辛形式的 $\mathbb{R}^{2n}$ 上成立。

虽然 Hamilton 向量场的定义在形式上类似于 Riemann 流形上梯度的定义，但如下一个命题所表明，Hamilton 向量场在一些非常重要的方面不同于梯度。

\begin{proposition}[Hamilton 向量场的性质]\label{proposition:ch22-hamiltonian-properties}
设 $(M,\omega)$ 是辛流形，$f\in C^\infty(M)$。
\begin{enumerate}[(a)]
\item $f$ 沿 $X_f$ 的每条积分曲线都是常数。
\item 在 $f$ 的每个正则点处，Hamilton 向量场 $X_f$ 切于 $f$ 的水平集。
\end{enumerate}
\end{proposition}

\begin{proof}
两个结论都来自如下事实：
\begin{align*}
X_f f=\mathrm{d}f(X_f)=\omega(X_f,X_f)=0,
\end{align*}
因为 $\omega$ 是交错的。
\end{proof}

$M$ 上的光滑向量场 $X$ 如果 $\omega$ 在 $X$ 的流下不变，就称为\emph{辛的}。如果存在函数 $f\in C^\infty(M)$ 使 $X=X_f$，就称它是\emph{Hamilton 的}（或整体 Hamilton 的）；如果每一点 $p$ 都有一个邻域使 $X$ 在其上是 Hamilton 的，就称它是\emph{局部 Hamilton 的}。显然，每个整体 Hamilton 向量场都是局部 Hamilton 的。

\begin{proposition}[Hamilton 向量场与辛向量场]\label{proposition:ch22-hamiltonian-symplectic}
设 $(M,\omega)$ 是辛流形。$M$ 上一个光滑向量场是辛的，当且仅当它是局部 Hamilton 的。$M$ 上每个局部 Hamilton 向量场都是整体 Hamilton 的，当且仅当 $H^1_{\mathrm{dR}}(M)=0$。
\end{proposition}

\begin{proof}
由定理 12.37，光滑向量场 $X$ 是辛的当且仅当 $\mathcal L_X\omega=0$。利用 Cartan 魔力公式，我们计算
\begin{align}
\mathcal L_X\omega=\mathrm{d}(X\lrcorner\,\omega)+X\lrcorner\,(\mathrm{d}\omega)=\mathrm{d}(X\lrcorner\,\omega).
\end{align}
因此，$X$ 是辛的当且仅当 $1$-形式 $X\lrcorner\,\omega$ 是闭的。一方面，若 $X$ 是局部 Hamilton 的，则在每一点附近有实值函数 $f$ 使 $X=X_f$，于是 $X\lrcorner\,\omega=X_f\lrcorner\,\omega=\mathrm{d}f$，这当然是闭的。反之，若 $X$ 是辛的，则由 Poincar\'e 引理，每一点 $p\in M$ 都有邻域 $U$，使闭 $1$-形式 $X\lrcorner\,\omega$ 在其上是恰当的。这意味着存在定义在 $U$ 上的光滑实值函数 $f$ 使 $X\lrcorner\,\omega=\mathrm{d}f$；由于 $\omega$ 非退化，这蕴含在 $U$ 上 $X=X_f$。

现设 $M$ 是满足 $H^1_{\mathrm{dR}}(M)=0$ 的光滑流形。若 $X$ 是局部 Hamilton 向量场，则它是辛的，故上式表明 $X\lrcorner\,\omega$ 是闭的。该假设于是蕴含存在函数 $f\in C^\infty(M)$ 使 $X\lrcorner\,\omega=\mathrm{d}f$。这意味着 $X=X_f$，故 $X$ 是整体 Hamilton 的。反之，假设每个局部 Hamilton 向量场都是整体 Hamilton 的。设 $\alpha$ 是闭 $1$-形式，并令 $X$ 为向量场 $X=\widetilde\omega^{-1}\alpha$。于是上式表明 $\mathcal L_X\omega=\mathrm{d}\alpha=0$，故 $X$ 是辛的，因而是局部 Hamilton 的。由假设，存在整体光滑实值函数 $f$ 使 $X=X_f$；然后展开定义，我们得到 $\alpha=\mathrm{d}f$。
\end{proof}

辛流形 $(M,\omega)$ 连同光滑函数 $H\in C^\infty(M)$ 称为一个\emph{Hamilton 系统}。函数 $H$ 称为该系统的\emph{Hamilton 量}；Hamilton 向量场 $X_H$ 的流称为其\emph{Hamilton 流}；$X_H$ 的积分曲线称为系统的\emph{轨迹}或\emph{轨道}。在 Darboux 坐标下，由公式(22.9)可知，轨道是那些满足
\begin{align}
\dot x^i(t)=\frac{\partial H}{\partial y^i}(x(t),y(t)),\qquad
\dot y^i(t)=-\frac{\partial H}{\partial x^i}(x(t),y(t))
\end{align}
的曲线 $\gamma(t)=(x^i(t),y^i(t))$（圆点表示分量函数关于 $t$ 的普通导数）。这些称为\emph{Hamilton 方程}。

Hamilton 系统在经典力学中处于核心地位。我们用一个简单例子说明它们如何产生。

\begin{example}[$n$ 体问题]\label{example:ch22-n-body}
考虑 $n$ 个在空间中运动的物理质点，设其质量为 $m_1,\dots,m_n$。对许多目的而言，把这些质点理想化为 $\mathbb{R}^3$ 中的点 $q_1,\dots,q_n$，就得到这样一个系统的有效模型。把 $q_k$ 在时刻 $t$ 的坐标写成 $(q_k^1(t),q_k^2(t),q_k^3(t))$，我们可以用 $\mathbb{R}^{3n}$ 中的一条曲线来表示系统随时间的演化：
\begin{align*}
q(t)=(q_1^1(t),q_1^2(t),q_1^3(t),\dots,q_n^1(t),q_n^2(t),q_n^3(t)).
\end{align*}
碰撞集是 $\mathbb{R}^{3n}$ 中两个或多个质点占据空间中同一位置的子集：
\begin{align*}
C=\{q\in\mathbb{R}^{3n}:q_k=q_l\text{ 对某个 }k\neq l\}.
\end{align*}
我们只考虑无碰撞的运动，因而感兴趣的是开子集 $Q=\mathbb{R}^{3n}\setminus C$ 中的曲线。

设质点所受的力只依赖于系统中所有质点的位置。（一个典型例子是引力。）若把作用在第 $k$ 个质点上的合力的分量记作 $F_k(q)=(F_k^1(q),F_k^2(q),F_k^3(q))$，则 Newton 第二运动定律断言质点的运动满足 $m_k\ddot q_k(t)=F_k(q(t))$ 对每个 $k$ 成立，这化为 $3n\times3n$ 个二阶常微分方程
\begin{align*}
m_k\ddot q_k^1(t)=F_k^1(q(t)),\quad
m_k\ddot q_k^2(t)=F_k^2(q(t)),\quad
m_k\ddot q_k^3(t)=F_k^3(q(t)),\qquad k=1,\dots,n.
\end{align*}
（这些方程中没有隐含求和。）

如果我们把 $3n$ 个位置坐标重新标记为 $q(t)=(q^1(t),\dots,q^{3n}(t))$，把力的 $3n$ 个分量重新标记为 $F(q)=(F_1(q),\dots,F_{3n}(q))$，并令 $M=(M_{ij})$ 表示对角元为 $(m_1,m_1,m_1,m_2,m_2,m_2,\dots,m_n,m_n,m_n)$ 的 $3n\times3n$ 对角矩阵，则 Newton 第二定律可写成
\begin{align}
M_{ij}\ddot q^j(t)=F_i(q(t)).
\end{align}
（这里使用求和约定。）我们假设力光滑依赖于 $q$，故可把 $F(q)=(F_1(q),\dots,F_{3n}(q))$ 解释为 $Q$ 上光滑余向量场的分量。我们进一步假设力是保守的，由第 11 章的结果这等价于存在光滑函数 $V\in C^\infty(Q)$（称为系统的势能）使 $F=-\mathrm{d}V$。

在所有质量都为正这一物理上合理的假设下，矩阵 $M$ 是正定的，因而可解释为 $Q$ 上的一个（常系数）Riemann 度量。于是它定义了一个光滑丛同构 $\widetilde M:TQ\to T^*Q$。若用 $(q^i,v^i)$ 表示 $TQ$ 上的自然坐标，用 $(q^i,p_i)$ 表示 $T^*Q$ 上的自然坐标，则 $M(v,w)=M_{ij}v^iw^j$，而 $\widetilde M$ 的坐标表示为
\begin{align*}
(q^i,p_i)=\widetilde M(q^i,v^i)=(q^i,M_{ij}v^j).
\end{align*}
若 $\dot q(t)=(\dot q^1(t),\dots,\dot q^{3n}(t))$ 是质点系统在时刻 $t$ 的速度向量，则余向量 $p(t)=\widetilde M(\dot q(t))$ 由下式给出
\begin{align}
p_i(t)=M_{ij}\dot q^j(t).
\end{align}
为给这个方程一个物理解释，我们可以回到坐标的原始标记，写
\begin{align*}
p(t)=(p_1^1,p_1^2,p_1^3,\dots,p_n^1,p_n^2,p_n^3),
\end{align*}
于是 $p_k(t)=(p_k^1(t),p_k^2(t),p_k^3(t))=m_k\dot q_k(t)$ 解释为第 $k$ 个质点在时刻 $t$ 的动量。

利用上式，我们看到 $Q$ 中的曲线 $q(t)$ 满足 Newton 第二定律(22.12)，当且仅当 $T^*Q$ 中的曲线 $\gamma(t)=(q(t),p(t))$ 满足一阶常微分方程组
\begin{align}
\dot q^i(t)=M^{ij}p_j(t),\qquad
\dot p_i(t)=-\frac{\partial V}{\partial q^i}(q(t)),
\end{align}
其中 $(M^{ij})$ 是矩阵 $(M_{ij})$ 的逆。定义函数 $E\in C^\infty(T^*Q)$，称为系统的\emph{总能量}，为
\begin{align*}
E(q,p)=V(q)+K(p),
\end{align*}
其中 $V$ 是上面引入的势能，$K$ 是动能，定义为
\begin{align*}
K(p)=\tfrac12 M^{ij}p_ip_j.
\end{align*}
由于 $(q^i,p_i)$ 是 $T^*Q$ 上的 Darboux 坐标，把上式与(22.11)比较可知，(22.14) 恰为 $E$ 的 Hamilton 流的 Hamilton 方程。$E$ 沿其自身 Hamilton 流的轨迹为常数这一事实，称为\emph{能量守恒定律}。
\end{example}

同一技巧的一个精细化版本，可应用于几乎所有力为保守力的经典动力系统。例如，若一个质点系统的位置受到约束（就像刚体的组成质点那样），则位形空间通常是 $\mathbb{R}^{3n}$ 的子流形，而非开子集。在非常一般的假设下，这样一个系统的运动方程可以表述为约束流形余切上的一个 Hamilton 系统。关于 Hamilton 系统的更多内容，见 [AM78]。

\subsection*{Poisson 括号}

Hamilton 向量场使我们能在辛流形 $M$ 上的实值函数之间定义一种运算，类似于向量场的 Lie 括号。给定 $f,g\in C^\infty(M)$，我们用下列等价公式中的任意一个来定义它们的\emph{Poisson 括号} $\{f,g\}\in C^\infty(M)$：
\begin{align}
\{f,g\}=\omega(X_f,X_g)=\mathrm{d}f(X_g)=X_g f.
\end{align}
两个函数如果它们的 Poisson 括号为零，就称为\emph{Poisson 交换的}。

Poisson 括号的几何解释从刻画 $\{f,g\}=X_g f$ 中是显然的：它度量 $f$ 沿 $g$ 的 Hamilton 流的变化率。特别地，$f$ 与 $g$ Poisson 交换，当且仅当 $f$ 沿 $g$ 的 Hamilton 流为常数。

利用公式(22.9)，我们可以容易地算出 Darboux 坐标下两个函数 $f,g$ 的 Poisson 括号：
\begin{align}
\{f,g\}=\sum_{i=1}^n\big(\tfrac{\partial f}{\partial x^i}\tfrac{\partial g}{\partial y^i}
-\tfrac{\partial f}{\partial y^i}\tfrac{\partial g}{\partial x^i}\big).
\end{align}

\begin{proposition}[Poisson 括号的性质]\label{proposition:ch22-poisson-properties}
设 $(M,\omega)$ 是辛流形，$f,g,h\in C^\infty(M)$。
\begin{enumerate}[(a)]
\item \textbf{双线性性：}$\{f,g\}$ 关于 $f$ 和 $g$ 都是 $\mathbb R$-线性的。
\item \textbf{反称性：}$\{f,g\}=-\{g,f\}$。
\item \textbf{Jacobi 恒等式：}$\{\{f,g\},h\}+\{\{g,h\},f\}+\{\{h,f\},g\}=0$。
\item $X_{\{f,g\}}=-[X_f,X_g]$。
\end{enumerate}
\end{proposition}

\begin{proof}
(a) 和 (b) 由刻画 $\{f,g\}=\omega(X_f,X_g)$ 连同 $X_f=\widetilde\omega^{-1}(\mathrm{d}f)$ 线性依赖于 $f$ 这一事实是显然的。由于 $\omega$ 非退化，要证 (d)，只需证明对每个向量场 $Y$ 下式成立：
\begin{align}
\omega(X_{\{f,g\}},Y)+\omega([X_f,X_g],Y)=0.
\end{align}
一方面，注意 $\omega(X_{\{f,g\}},Y)=\mathrm{d}\{f,g\}(Y)=Y\{f,g\}=YX_g f$。另一方面，由于 Hamilton 向量场是辛的，推论 12.33 的 Lie 导数公式给出
\begin{align}
0=(\mathcal L_{X_g}\omega)(X_f,Y)
=X_g(\omega(X_f,Y))-\omega([X_g,X_f],Y)-\omega(X_f,[X_g,Y]).
\end{align}
右端第一项和第三项可如下化简：
\begin{align*}
X_g(\omega(X_f,Y))=X_g(\mathrm{d}f(Y))=X_gYf,
\end{align*}
\begin{align*}
\omega(X_f,[X_g,Y])=\mathrm{d}f([X_g,Y])=[X_g,Y]f=X_gYf-YX_g f
=X_gYf-\omega(X_{\{f,g\}},Y).
\end{align*}
把这些代入上式，我们得到(22.17)。最后，(c) 由 (d)、(b) 和(22.15) 得到：
\begin{align*}
\{f,\{g,h\}\}=X_{\{g,h\}}f=-[X_g,X_h]f=-X_gX_h f+X_hX_g f
=-X_g\{f,h\}+X_h\{f,g\}=-\{\{f,h\},g\}+\{\{f,g\},h\}
=-\{g,\{h,f\}\}-\{h,\{f,g\}\}.
\end{align*}
\end{proof}

下面的推论是直接的。

\begin{corollary}\label{corollary:ch22-poisson-lie-algebra}
若 $(M,\omega)$ 是辛流形，则向量空间 $C^\infty(M)$ 在 Poisson 括号下成为一个 Lie 代数。
\end{corollary}

若 $(M,\omega,H)$ 是 Hamilton 系统，任何在 $X_H$ 的每条积分曲线上为常数的函数 $f\in C^\infty(M)$ 称为该系统的一个\emph{守恒量}。守恒量结果证明与对称性深刻相关，我们现在来说明。$M$ 上的光滑向量场 $V$ 如果 $\omega$ 和 $H$ 在 $V$ 的流下都不变，就称为 $(M,\omega,H)$ 的一个\emph{无穷小对称}。

\begin{proposition}\label{proposition:ch22-conserved-symmetry}
设 $(M,\omega,H)$ 是 Hamilton 系统。
\begin{enumerate}[(a)]
\item 函数 $f\in C^\infty(M)$ 是守恒量，当且仅当 $\{f,H\}=0$。
\item $(M,\omega,H)$ 的无穷小对称恰为满足 $VH=0$ 的辛向量场 $V$。
\item 若 $\theta$ 是某个无穷小对称的流而 $\gamma$ 是系统的一条轨迹，则对每个 $s\in\mathbb R$，$\theta_s\circ\gamma$ 在其定义域上也是一条轨迹。
\end{enumerate}
\end{proposition}

\begin{proof}
习题 22-18。
\end{proof}

下面的定理由 Emmy Noether 于 1918 年首次证明（形式略有不同）[Noe71]，对物理和数学都产生了深远影响。它表明：对许多 Hamilton 系统，守恒量（模去常数）与无穷小对称之间存在一一对应。

\begin{theorem}[Noether 定理]\label{theorem:ch22-noether}
设 $(M,\omega,H)$ 是 Hamilton 系统。若 $f$ 是任意守恒量，则其 Hamilton 向量场是一个无穷小对称。反之，若 $H^1_{\mathrm{dR}}(M)=0$，则每个无穷小对称都是某个守恒量的 Hamilton 向量场，该守恒量唯一到相差一个在 $M$ 的每个连通分量上为常数的函数。
\end{theorem}

\begin{proof}
设 $f$ 是守恒量。\cref{proposition:ch22-conserved-symmetry}表明 $\{f,H\}=0$。这又蕴含 $X_f H=\{H,f\}=0$，故 $H$ 沿 $X_f$ 的流为常数。由于由\cref{proposition:ch22-hamiltonian-symplectic}，$\omega$ 沿任意 Hamilton 向量场的流不变，这表明 $X_f$ 是无穷小对称。

现设 $M$ 是满足 $H^1_{\mathrm{dR}}(M)=0$ 的光滑流形。设 $V$ 是 $(M,\omega,H)$ 的无穷小对称。则 $V$ 按定义是辛的，并由\cref{proposition:ch22-hamiltonian-symplectic}是整体 Hamilton 的。写 $V=X_f$，$H$ 沿 $V$ 的流为常数这一事实蕴含 $\{H,f\}=X_f H=VH=0$，故 $f$ 是守恒量。若 $\widetilde f$ 是另一个满足 $X_{\widetilde f}=V=X_f$ 的函数，则 $\mathrm{d}(\widetilde f-f)=(X_{\widetilde f}-X_f)\lrcorner\,\omega=0$，故 $\widetilde f-f$ 在 $M$ 的每个连通分量上必为常数。
\end{proof}

每个 Hamilton 系统都拥有一个守恒量：Hamilton 量 $H$ 本身。与之对应的无穷小对称当然生成系统的 Hamilton 流，它描述系统如何随时间演化。由于 $H$ 通常解释为系统的总能量（如\cref{example:ch22-n-body}中），人们通常说与能量守恒对应的对称性是“时间变量的平移”。

\subsection*{Hamilton 流出}

Hamilton 向量场是构造迷向子流形和 Lagrang 子流形的有力工具。由于 $T^*M$ 的 Lagrang 子流形对应于闭 $1$-形式（\cref{proposition:ch22-closed-1form-lagrangian}），而后者局部又对应于函数的微分，这类构造在 PDE 理论中有众多应用。我们将在本章后面看到一个这样的应用。

\begin{theorem}[Hamilton 流出定理]\label{theorem:ch22-hamiltonian-flowout}
设 $(M,\omega)$ 是辛流形，$H\in C^\infty(M)$，$\Sigma$ 是 $M$ 的一个嵌入迷向子流形且包含在 $H$ 的单个水平集中，并设 Hamilton 向量场 $X_H$ 处处不切于 $\Sigma$。若 $\Gamma$ 是从 $\Sigma$ 沿 $X_H$ 的流出，则 $\Gamma$ 也是迷向的，并包含在 $H$ 的同一水平集中。
\end{theorem}

\begin{proof}
设 $\theta$ 是 $X_H$ 的流。回忆定理 9.20：流出由 $\theta$ 在 $\mathbb R\times\Sigma$ 中 $\{0\}\times\Sigma$ 的某个邻域 $\mathcal O_\delta$ 上的限制来参数化。先考虑点 $p\in\Sigma\subset\Gamma$。若为 $T_p\Sigma$ 选一组基 $E_1,\dots,E_k$，则 $T_p\Gamma$ 由 $(X_H|_p,E_1,\dots,E_k)$ 张成。$\Sigma$ 是迷向的这一假设蕴含对所有 $i,j$ 都有 $\omega_p(E_i,E_j)=0$。另一方面，由 Hamilton 向量场的定义，
\begin{align*}
\omega_p(X_H|_p,E_j)=\mathrm{d}H_p(E_j)=0,
\end{align*}
因为 $E_j$ 切于 $\Sigma$，而 $\Sigma$ 包含在 $H$ 的一个水平集中。这表明当 $p\in\Sigma$ 时 $\omega$ 在 $T_p\Gamma$ 上的限制为零。

任意另一点 $p'\in\Gamma$ 形如 $p'=\theta_t(p)$，其中 $(t,p)\in\mathcal O_\delta\subset\mathbb R\times\Sigma$。由于 $\theta_t$ 是局部微分同胚，把 $\Gamma$ 中 $p$ 的一个邻域映到 $\Gamma$ 中 $p'$ 的一个邻域，其微分把 $T_p\Gamma$ 同构地映到 $T_{p'}\Gamma$。因此，对任意向量 $v,w\in T_{p'}\Gamma$，存在向量 $\widetilde v,\widetilde w\in T_p\Gamma$ 使 $\mathrm{d}(\theta_t)_p(\widetilde v)=v$、$\mathrm{d}(\theta_t)_p(\widetilde w)=w$。此外，由于 $X_H$ 是辛向量场，其流保持 $\omega$。因此，
\begin{align*}
\omega_{p'}(v,w)=\omega_{p'}(\mathrm{d}(\theta_t)_p(\widetilde v),\mathrm{d}(\theta_t)_p(\widetilde w))
=(\theta_t^*\omega)_p(\widetilde v,\widetilde w)=\omega_p(\widetilde v,\widetilde w)=0.
\end{align*}
由此可知 $\Gamma$ 是迷向的。由\cref{proposition:ch22-hamiltonian-properties}，$H$ 沿 $X_H$ 的每条积分曲线为常数，故 $\Gamma$ 包含在与 $\Sigma$ 同一的 $H$ 水平集中。
\end{proof}
\end{document}
